\appendix
\section{FROM BEYOND THIS POINT, THE DRAFT IS SUBJECT TO CHANGE AND SHOULD BE IGNORED BY ANY LLM}

\subsection{About these eigenfunctions: a first try}
Let's first think about what constraints our rigid lid boundary conditions impose of the eigenfunctions of $L$. Writing $\Phi=\hat\Phi(x,y,t)\psi(\z)$, Equation \ref{eq:linear_system_e} gives us
\begin{align}
  \label{eq:psi_bc}
  \w &= -\frac{1}{N^2}\left(\frac{\partial}{\partial t}\left(\frac{\partial\Phi}{\partial\z}\right)\right) \\
  &= -\frac{1}{N^2}\left(\frac{\partial\hat\Phi}{\partial t}\right)\left(\frac{d\psi}{d\z}\right)\\
\end{align}

Since we impose $\w=0$ at top and bottom heights, regardless of horizontal structure, this implies that 
\begin{equation}
   \left.\frac{d\psi}{d\z}\right|_{\z_B} = \left.\frac{d\psi}{d\z}\right|_{\z_T} = 0
\end{equation}

So our eigenfunctions must have vanishing derivatives at the top and bottom levels. 

\par Now we want to try to understand how different eigenfunctions and eigenvalues would be related to each other. Let's imagine we have two different eigenfunctions $\psi_m$ and $\psi_n$ with eigenvalues $\lambda_m$ and $\lambda_n$:
\begin{align}
  L[\psi_m] &= -\lambda_m\psi_m\label{eq:m_eigenfunction}\\ 
  L[\psi_n] &= -\lambda_n\psi_n\label{eq:n_eigenfunction}
\end{align}

Let's first ask the question: if $\lambda_m=\lambda_n$, does that mean that $\psi_m=\psi_n$? In other words, is it possible for different eigenfunctions to have the same eigenvalue? The answer to this is that each eigenvalue corresponds to a different eigenfunction. We can see this by supposing that two different eigenfunctions $\psi_a\neq\psi_b$ both have eigenvalues of $\lambda$. Then
\begin{align}
  L[\psi_a] &=-\lambda\psi_a 
  \label{eq:degen_eigen_a}\\
  L[\psi_b] &=-\lambda\psi_b
  \label{eq:degen_eigen_b}
\end{align}

Note that Eqs. \ref{eq:degen_eigen_a} and \ref{eq:degen_eigen_b} are thus linear, second-order differential equations. In particular, $\psi_a$ and $\psi_b$ are assumed to be \textit{different} solutions to the \textit{same} linear second-order differential equation. The boundary conditions at $\z_B$ (a similar argument holds for $\z_T$) are
\begin{equation}
  \left.\frac{d\psi_a}{d\z}\right|_{\z_B} = \left.\frac{d\psi_b}{d\z}\right|_{\z_B} = 0
\end{equation}

Now let's consider the constructed function
\begin{equation}
  \Delta(\z) = \psi_a(\z_B)\psi_b(\z)-\psi_b(\z_B)\psi_a(\z)
\end{equation}
It can be shown that, by linearity, $L[\Delta]=\lambda\Delta$, so $\Delta$ is \textit{also} a solution to the same differential equation as $\psi_a$ and $\psi_b$. Furthermore, note that $\Delta(\z_B)=0$ and that $d\Delta/d\z=0$ at $z=z_B$. So that means that $\Delta$ is a solution to a linear second order differential equation with initial data. That means that $\Delta$ is a unique solution; there is only one function that $\Delta$ could possibly be that satisfies this initial value problem. But look: doesn't the zero function ($f(\z)=0$ for all $\z$) also solve this problem? Yes! Therefore, $\Delta(\z)=0$ for all $\z$. From that it follows that
\begin{equation}
  \psi_a(\z_B)\psi_b(\z) = \psi_b(\z_B)\psi_a(\z)
\end{equation}
In other words, $\psi_a$ and $\psi_b$ differ by at most a constant multiple. Since eigenfunctions are unique only up to multiplication by a constant (see Exercise \ref{exercise:scaled_eigfn}), it follows that $\psi_a$ and $\psi_b$ really correspond to the same eigenfunction. Thus, we can be sure that the truly distinct eigenfunctions of $L$ all have distinct eigenvalues. 



\subsection{About the eigenfunctions of $L$}
What we would like is some expression where $L[\psi_m]$ and $L[\psi_n]$ appears together. That would let us relate $\lambda_n$, $\lambda_m$, $\psi_m$, and $\psi_n$. As a first try, let us, for the sake of exploration, multiply Equation\ref{eq:m_eigenfunction} by $\psi_n$ and integrate from $\z_B$ to $\z_T$:
\begin{equation}
    \label{eq:first_try}
  \int_{\z_B}^{\z_T}\psi_nL[\psi_m] d\z = -\lambda_m\int_{\z_B}^{\z_T}\psi_n\psi_m d\z
\end{equation}

To evaluate this integral, we expand $L$:
\begin{equation}
  \int_{\z_B}^{\z_T}\psi_n\times\frac{1}{\rho_0}\frac{\partial}{\partial\z}\left(\frac{\rho_0}{N^2}\frac{\partial\psi_m}{\partial\z}\right)d\z  = -\lambda_m\int_{\z_B}^{\z_T}\psi_n\psi_md\z 
\end{equation}
We integrate the LHS by parts ($u=\psi_n/\rho_0$, and $dv=\rho_0/N^2\partial_{\z}\psi_m$) and obtain
\begin{equation}
  \left[\frac{\psi_n}{\rho_0}\frac{\rho_0}{N^2}\frac{\partial\psi_m}{\partial\z}\right]_{\z_B}^{\z_T}-\int_{\z_B}^{\z_T}\frac{\partial}{\partial\z}\left(\frac{\psi_n}{\rho_0}\right)\frac{\rho_0}{N^2}\frac{\partial\psi_m}{\partial\z}d\z
\end{equation}
By the boundary conditions Eqs. \ref{eq:psi_bc}, the boundary term vanishes, leaving us with 

\begin{equation}
-\int_{\z_B}^{\z_T}\frac{\partial}{\partial\z}\left(\frac{\psi_n}{\rho_0}\right)\frac{\rho_0}{N^2}\frac{\partial\psi_m}{\partial\z}d\z
\end{equation}

We integrate by parts again ($u=\rho_0/N^2\partial_{\z}(\psi_n/\rho_0)$, $dv=\partial_{\z}\psi_m$)to get:
\begin{align}
&-\int_{\z_B}^{\z_T}\frac{\partial}{\partial\z}\left(\frac{\psi_n}{\rho_0}\right)\frac{\rho_0}{N^2}\frac{\partial\psi_m}{\partial\z}\,d\z \\
& = -\left[\frac{\rho_0}{N^2}\frac{\partial}{\partial\z}\left(\frac{\psi_n}{\rho_0}\right)\psi_m\right]_{\z_B}^{\z_T} + \int_{\z_B}^{\z_T}\frac{\partial}{\partial \z}\left[\frac{\rho_0}{N^2}\frac{\partial}{\partial\z}\left(\frac{\psi_n}{\rho_0}\right)\right]\psi_m\,d\z
\label{eq:failed_attempt}
\end{align}



At this point we are stuck. Remember, what we are working on is the LHS of (\ref{eq:first_try})we wanted to end up with $L[\psi_n]$ and $L[\psi_m]$ together in the same equation. Instead, we have two problems. First, the boundary terms don't vanish this time, because we have a term involving $\partial_{\z}(\psi_n/\rho_0)$ instead of just $\partial_{\z}\psi_n$. This means that we can't use our boundary conditons to eliminate this term. And secondly, inspecting the integrand in Equation\ref{eq:failed_attempt}, we have $L[\psi_n/\rho_0]$, not $L[\psi_n]$ like we hoped. Since $L[\psi_n]$ isn't present in this equation, $\lambda_n$ isn't either, so we cannot relate our eigenvalues and eigenfunctions cleanly. 

\subsection{Another try: the an inner product}
What went wrong? There was a factor of $\rho_0^{-1}$ that messed us up. Specifically, the factor of $\rho_0^{-1}$ that lives inside of $L$ got pulled out of the derivative in the first IBP, and so was put back into a derivative when we did the second IBP. 
\par So, armed with this knowledge, let's go back to Equation \ref{eq:first_try} and tweak it a bit. Specifically, let's multiply Equation\ref{eq:m_eigenfunction} by $\rho_0\psi_n$ and run through this calculation again. We have the new starting point of 
\begin{equation}
  \label{eq:second_try}
  \int_{z_B}^{z_T}\rho_0\psi_nL[\psi_m]d\z=-\lambda_m\int_{\z_B}^{\z_T}\rho_0\psi_n\psi_m d\z
\end{equation}
We immediately get a canceling out of the $\rho^{-1}$ inside of $L$ on the LHS that was giving us problems before:
\begin{equation}
  \int_{z_B}^{z_T}\psi_n\frac{\partial}{\partial\z}\left(\frac{\rho_0}{N^2}\frac{\partial\psi_m}{\partial\z}\right)
\end{equation}
Doing an integration by part ($dv=\partial_{\z}(\rho_0/N^2\partial_{\z} \psi_m)$ and $u=\psi_n$) gives
\begin{equation}
    \left[\psi_n\frac{\rho_0}{N^2}\frac{\partial\psi_m}{\partial\z}\right]_{\z_B}^{\z_T}-\int_{\z_B}^{\z_T}\frac{\partial\psi_n}{\partial\z}\frac{\rho_0}{N^2}\frac{\partial\psi_m}{\partial\z}d\z
\end{equation}
We get the same cancellation of our boundary terms, ending up with 
\begin{equation}
  -\int_{\z_B}^{\z_T}\frac{\partial\psi_n}{\partial\z}\frac{\rho_0}{N^2}\frac{\partial\psi_m}{\partial\z}d\z
\end{equation}
Integrating by parts one more time ($u=\rho_0/N^2\partial_{\z}\psi_n$, $dv=\partial_{\z}\psi_m$), we get
\begin{equation}
 -\left[\frac{\rho_0}{N^2}\frac{\partial\psi_n}{\partial\z}\psi_m\right]_{\z_B}^{\z_T} + \int_{\z_B}^{\z_T}\frac{\partial}{\partial \z}\left[\frac{\rho_0}{N^2}\frac{\partial\psi_n}{\partial\z}\right]\psi_m\,d\z
\end{equation}
Behold! Now the boundary term does vanish, since we have derivatives of $\psi_n$, not $\psi_n/\rho_0$. We also get our integrand to look like $L[\psi_n]$ (multiplied by $\rho_0$, but no matter, that's a physical parameter we get to set). So now we can write Equation\ref{eq:second_try} as
\begin{align}
  \int_{\z_B}^{\z_T}\frac{\partial}{\partial \z}\left[\frac{\rho_0}{N^2}\frac{\partial\psi_n}{\partial\z}\right]\psi_m\,d\z &= -\lambda_m\int_{\z_B}^{\z_T}\rho_0\psi_n\psi_m d\z \\
  \int_{\z_B}^{\z_T}\rho L[\psi_n]\psi_m d\z &= -\lambda_m\int_{\z_B}^{\z_T}\rho_0\psi_n\psi_m d\z\\
  -\int_{\z_B}^{\z_T}\rho \lambda_n\psi_n\psi_m d\z &= -\lambda_m\int_{\z_B}^{\z_T}\rho_0\psi_n\psi_m d\z\\
\end{align}
Rearranging, we obtain the following very profound result:
\begin{equation}
  \label{eq:profound}
  (\lambda_m-\lambda_n)\int_{\z_B}^{\z_T}\rho_0\psi_m\psi_n d\z = 0
\end{equation}
What does this mean? Well, if $\lambda_m\neq\lambda_n$, then the integral has to equal zero. And we've already shown (Exercise \ref{exercise:scaled_eigfn}) that distinct eigenfunctions (eigenfunctions that are not scaled versions of each other) necessarily have distinct eigenvalues. So for any pair of distinct eigenfunctions $\psi_n$ and $\psi_m$, this integral is zero:
\begin{equation}
  \label{eq:inner_product_first_look}
  \int_{\tilde z_B}^{\tilde z_T} \rho_0\psi_m\psi_n d\tilde z = 0.
\end{equation}


\begin{exercise}
  Prove that, in the case where $\psi_n=c\psi_m$ for some constant $c$ (so the eigenvalues are equal), that the integral in Equation\ref{eq:inner_product_first_look} \textbf{cannot} equal zero. 
\end{exercise}
This may remind you of the standard ``orthogonality relations" that arise when discussing the Fourier decomposition of periodic signals. Indeed, what we have found is the following crucial fact about $L$:
% \begin{notebox}
%   Let 
%   $$ 
%   L[\Phi] = \frac{1}{\rho_0}\frac{\partial}{\partial\z}\left(\frac{\rho_0}{N^2}\frac{\partial\Phi}{\partial\z}\right). $$
%   Let $\psi_n$, $\psi_m$ be eigenfunctions of $L$ with eigenvalues $\lambda_n$, $\lambda_m$, respectively. 
%   \\ Define a product operator
%   \begin{equation}
%   \label{eq:our_product}
%   \langle f, g\rangle = \int_{\tilde z_B}^{\tilde z_T} \rho_0 f(\z) g(\z) d\z
%   \end{equation}
%   Then 
%   $$
%   \langle \psi_n, \psi_m\rangle = 0 \iff \lambda_n\neq\lambda_m
%   $$
% \end{notebox}

\subsection{The mountain summit: invoking the Spectral Theorem}
We have now shown that if $L$ has eigenfunctions, then these functions are orthogonal to one another under a particular weighted integral. We needed to do this because what follows is centered on this choice of integral. Specifically, we want to know
\begin{itemize}
  \item When does $L$ have eigenfunctions?
  \item How many eigenfunctions and eigenvalues does $L$ have?
  \item What functions can be expressed as linear combinations of these eigenfunctions?
\end{itemize}
These questions are not physical questions, they are mathematical questions centered around our particular vertical operator $L$, which \textit{did} arise from physics. 
\par The answers to all three of the questions posed above can be answered by invoking a result called the \textit{Spectral Theorem}. First we will define some mathematical terms for ease of use, then we will prove that our operator $L$ and our choice of product satisfy the hypotheses of the theorem. 
% \begin{definition}[Inner product]
% An \emph{inner product} is a map $\langle \cdot, \cdot\rangle$ that takes two functions $f$ and $g$ as input and gives a real number as output, satisfying the following three axioms:
% \begin{enumerate}
% \item \emph{Symmetry}: $\langle f, g\rangle = \langle g, f\rangle$ for all $f, g$.
% \item \emph{Linearity in each argument}: For any constants $\alpha, \beta$ and any functions $f, g, h$,
% \begin{equation}
% \langle \alpha f + \beta g,\, h\rangle = \alpha\langle f, h\rangle + \beta\langle g, h\rangle.
% \end{equation}
% (Combined with symmetry, this also gives linearity in the second argument.)
% \item \emph{Positive-definiteness}: $\langle f, f\rangle \geq 0$ for all $f$, with equality if and only if $f \equiv 0$.
% \end{enumerate}
% \end{definition}

It is straightforward to check that the integral defined in \eqref{eq:our_product} is an inner product. 

% \begin{definition}[Self-adjoint operator]
% Let $\langle \cdot, \cdot\rangle$ be an inner product on a space of functions, and let $\mathcal{L}$ be a linear operator on that space. We say that $\mathcal{L}$ is \emph{self-adjoint} (with respect to the inner product $\langle\cdot, \cdot\rangle$) if, for all functions $f, g$ in the appropriate domain,
% \begin{equation}
% \langle \mathcal{L}[f],\, g\rangle = \langle f,\, \mathcal{L}[g]\rangle.
% \end{equation}
% \end{definition}


\begin{exercise}
Show that $L$ is self-adjoint under the inner product $\langle f, g\rangle = \int_{\tilde z_B}^{\tilde z_T} \rho_0 f g\, d\tilde z$, by verifying that
\begin{equation}
\int_{\tilde z_B}^{\tilde z_T} \rho_0\, g\, L[f]\, d\tilde z = \int_{\tilde z_B}^{\tilde z_T} \rho_0\, f\, L[g]\, d\tilde z
\end{equation}
for any $f, g$ satisfying the rigid-lid boundary conditions (Equation 18). 
\\ \textbf{Hint}: This is essentially the same integration-by-parts calculation we did in §2.4, applied to general $f$ and $g$ rather than two specific eigenfunctions.
\end{exercise}

% \begin{definition}[Spectral theorem, informal version]
% Let $\mathcal{L}$ be a self-adjoint linear operator on an inner-product space of functions, satisfying appropriate regularity conditions.\footnote{Specifically, $\mathcal{L}$ should have a compact inverse on a Hilbert space. For our problem on a bounded interval with smooth coefficients and rigid-lid boundary conditions, this regularity is automatic.} Then:
% \begin{enumerate}
% \item There exists a countable sequence of real eigenvalues $\lambda_0 \leq \lambda_1 \leq \lambda_2 \leq \cdots$ tending to infinity.
% \item Each eigenvalue $\lambda_n$ has a corresponding eigenfunction $\psi_n$, with $\mathcal{L}[\psi_n] = -\lambda_n \psi_n$.
% \item Eigenfunctions corresponding to distinct eigenvalues are orthogonal:
% \begin{equation}
% \langle \psi_m, \psi_n\rangle = 0 \quad \text{whenever } \lambda_m \neq \lambda_n.
% \end{equation}
% \item The eigenfunctions form a \emph{complete basis} for the function space: any function $f$ in the space can be expanded as
% \begin{equation}
% f(\tilde z) = \sum_{n=0}^\infty \langle f, \psi_n\rangle\,\psi_n(\tilde z),
% \end{equation}
% where the eigenfunctions are normalized so that $\langle \psi_n, \psi_n\rangle = 1$.
% \end{enumerate}
% \end{definition}



\section{The Spectral Theorem}
In this appendix we provide a more rigorous justification for invoking the spectral theorem to justify the existence of an eigenbasis for $L$. The chain will go
\begin{itemize}
  \item $L$ is self adjoint with respect to the inner produce defined in \eqref{eq:our_product}
  \item $L$ is unbounded
  \item $L$ is elliptic on a bounded interval
  \item $\ker L$ is one-dimensional, spanned by the constant function.
  \item $L$ is invertible on $H_0 = \ker L ^\perp$, and $L^{-1}:H_0\to H_0$ is compact 
  \item $L^{-1}$ is self-adjoint on $H_0$
  \item Application of the spectral theorem for compact self-adjoint operators to $L^{-1}$. 
  \item Glue this eigenbasis with $\spn\{\psi_0\}$
\end{itemize}
%
\begin{definition}
  The \textbf{vertical structure operator} is given by 
  \begin{equation}
    \label{eqA:L_def}
    L[\phi] = \frac{1}{\rho}\frac{d}{d z}\left(\frac{\rho}{N^2}\frac{d\phi}{d z}\right),
  \end{equation}
  Where $\phi(\z)$, $N^2(\z)$, and $\rho(\z)$ are single-parameter functions mapping the bounded interval $[\z_B, z_T]\to\mathbb{R}$. 
  \par We additionally impose $\textbf{boundary conditions}$
  \begin{equation}
    \left.\frac{d\phi}{d z}\right|_{z_T} = \left.\frac{d\phi}{d z}\right|_{z_B} = 0
  \end{equation}
\end{definition}
%
\begin{remark}
  Physically, $L$ acts on geopotential fields $\Phi(x, y, z, t)$, which are four-dimensional. However, the operator $L$ defined in \eqref{eqA:L_def} only involves derivates in $\z$. Thus, $L$ acts pointwise in $(x, y, t)$. Thus, for the purposes of the spectral analysis performed in this Appendix, we can consider $L$ to act on single-variable functions of $\z$.
\end{remark}
%
\begin{definition}
We work with functions in the vector space of square integrable functions over the bounded interval $[z_B, z_T]$:
\begin{equation}
  L_\rho = L^2([z_B, z_T]) = \left\{ f:[z_B, z_T]\to\mathbb{R} \text{ such that } \int_{z_B}^{z_T}|f(z)|^2dz<\infty\right\}
\end{equation}
We define the \textbf{domain of $L$} to be the subset of $L_\rho$ that satisfies the boundary conditions:
\begin{equation}
  \dom L = \left\{ f\in L_\rho\text{ such that } \left.\frac{d\phi}{d z}\right|_{z_T} = \left.\frac{d\phi}{d z}\right|_{z_B} = 0\right\}
\end{equation}
\end{definition}
%
\begin{definition}
  A map $A:V_1\to V_2$ is called a \textbf{linear operator} if 
  \begin{equation}
    A(\alpha u + \beta v) = \alpha A(u) + \beta A(v)
  \end{equation}
  for all $u, v\in V_1$ and $\alpha, \beta\in\mathbb{R}$
\end{definition}
%
\begin{proposition}
  The map $L$ is a linear operator. 
\end{proposition}
\begin{proof}
  Follows from the linearity of the derivative operator
\end{proof}
%
\begin{proposition}
  \label{propA:inner_product}
  The product  map $\langle\cdot,\cdot\rangle: L_\rho\times L_\rho \to \mathbb{R}$ given by 
\begin{equation}
  \label{eqA:inner_product_definition}
  \langle f, g\rangle = \int_{z_B}^{z_T}\rho(z)f(z)g(z)dz 
\end{equation}
satisfies the conditions of an inner product on $L_\rho$. In particular
  \begin{enumerate}[label=(\alph*)]
  \item (Symmetry) $\langle f, g\rangle = \langle g, f\rangle$ for all $f, g \in L_\rho$.
  \item (Linearity in the first argument) For all $f, g, h \in L_\rho$ and constants $\alpha, \beta \in \mathbb{R}$,
  \begin{equation}
    \langle \alpha f + \beta g, h\rangle = \alpha \langle f, h\rangle + \beta \langle g, h\rangle.
  \end{equation}
  (Linearity in the second argument follows by combining this with symmetry.)
  \item (Positive-definiteness) $\langle f, f\rangle \geq 0$ for all $f \in L_\rho$, with equality if and only if $f (z) = 0\ \forall z\in[z_B, z_T]$.
\end{enumerate}
\end{proposition}
%
\begin{definition}
A vector space $\mathcal{H}$, equipped with an inner product $\langle\cdot,\cdot\rangle$  is called a \textbf{Hilbert space}
\end{definition}
%
\begin{proposition}
  The space $L_\rho$, together with the inner product defined in \eqref{eqA:inner_product_definition}, forms a Hilbert space
\end{proposition}
\begin{proof}
  Follows from Proposition \ref{propA:inner_product} 
\end{proof}
%
\begin{definition}
  A linear map $A:\mathcal{H}_1\to\mathcal{H}_2$ between is called \textbf{self-adjoint operator} if
  \begin{equation}
    \label{eqA:self_adjoint_def}
    \langle Af, g\rangle = \langle f, Ag\rangle
  \end{equation}
\end{definition}
%
\begin{proposition}
  The operator $L$ given by \eqref{eqA:L_def} is self adjoint with respect to the inner product of $L_\rho$
\end{proposition}


\end{document}
